\chapter{Introduction}\label{ch:introduction}

\section{Background and Motivation}

Large language models sometimes complete a factual prompt with a generic word while the correct answer sits a few places lower in the output ranking. For ``The location of Massachusetts Institute of Technology is in'', GPT-2 Small \cite{r1} may rank `` the'' first and `` Cambridge'' eighth. We call such cases \textbf{suppressed factual recall}; the term describes the output ranking only and makes no claim about what the model ``knows''. The usual remedies change the model: fine-tuning, or closed-form weight edits such as ROME \cite{r2} and MEMIT \cite{r3}, which are permanent and, applied repeatedly, first erase earlier edits gradually and then collapse the model \cite{r4}. Retrieval and in-context editing \cite{r34,r31} leave the weights alone but supply the fact as extra text.

Mechanistic interpretability offers a third route. A sparse autoencoder (SAE) \cite{r6,r7} re-expresses a layer's activations in a much wider, sparsely active basis whose directions, called \emph{features}, often match single concepts. A feature is a direction in the residual stream, so it can be turned down (\emph{muted}) or up (\emph{boosted}) during one forward pass and discarded afterwards. If this works, a wrong answer can be repaired \textbf{transiently}, without touching any weight, and the repair can be read as a short list of named features. This project asks whether it works, how it should be built, and what it really does.

\section{How the Project Grew}

The project, \textbf{FeatureScalpel} (early working name: Causal Feature Steering), started in July 2026 with a modest prototype and grew by trial and error. Every step was recorded in a numbered research journal (37 entries by October 2026), and this report follows that record in order. Figure~\ref{fig:stages} shows the whole path.

\begin{figure}[t]
\centering
\safegraphic[width=0.82\linewidth]{figures/h_stages}
\caption{The path of the project in twelve stages (J = research journal entries). Stages 1 to 11 are the main line (Chapters~\ref{ch:method} and~\ref{ch:evolution}); stage 12 is an applied side study (Chapter~\ref{ch:transfer}).}\label{fig:stages}
\end{figure}

Two things should be clear from the start. The \textbf{main line} (stages 1 to 11) is the project: a transient SAE edit of GPT-2 Small, how to choose the edit (by hand-built search, by a trained network, or by gradient descent), and what the edit does besides the one number it is tuned for. The \textbf{cross-model transfer study} (stage 12) is not a later version of the system. It is an \emph{applied side study}: once the edit, the controls and the benchmark existed, we used them on a new problem to see whether the interpretability and intervention work had fruits beyond the single model it was built for.

\section{Problem Statement and Research Questions}

\textit{Given a prompt whose true answer GPT-2 Small ranks below first, can a transient edit of SAE features during one forward pass bring the answer to the top, how should the edit be chosen, and what does reaching the top mean?} Five research questions organise the work.
\begin{enumerate}
  \item \textbf{RQ1 (what works).} Which way of choosing features and strengths brings a suppressed answer to rank 1, and which ideas fail?
  \item \textbf{RQ2 (limits).} What limits such an edit: the depth of the intervention, the safety rule, or the supply of features it can use?
  \item \textbf{RQ3 (choosing the edit).} Can a program choose the edit better than a hand-built search: a trained network, or gradient descent on the edit itself?
  \item \textbf{RQ4 (meaning of success).} Is rank 1 evidence that a fact was recovered? What else does the edit change: reworded prompts, neighbouring facts, unrelated prompts, generated text?
  \item \textbf{RQ5 (application).} Can the change an edit makes inside one model be carried into another model at inference time, without weight change, and what exactly is carried?
\end{enumerate}

\section{Objectives, Scope and Contributions}

The objectives are to build a transient edit that changes only the intended feature contribution; to document with numbers what each design decision bought and which ideas failed; to build a benchmark with held-out prompts and exact paired statistics; to measure side effects with controls; and to apply the tools to the transfer question with pass bars fixed before each run. The study uses GPT-2 Small with the SAE release \texttt{gpt2-small-res-jb} (24,576 features per layer) at layer 8 in all later studies, single-token facts from ROME's \texttt{known\_1000} and CounterFact \cite{r2}, and next-token measurements; the side study adds GPT-2 Medium (no SAE). Other models, comparison with ROME, MEMIT or in-context editing, and a full evaluation of generated text are outside the scope.

The contributions are: (1) a documented trial-and-error record from single-feature ablation to gradient descent (Chapter~\ref{ch:evolution}); (2) a transient edit with measured design choices (delta patching, causal ranking, a strict per-feature safety filter, pool refill, candidates from all prompt positions: 10 of 13 prompts at rank 1 against 4) and a batched implementation that cut model calls per prompt and layer from 186 to 6 with identical results; (3) an equal-budget result: on 300 held-out CounterFact prompts gradient descent reached rank 1 on 271, against 161 for a sweep given the same range and 94 for the original settings, with lower same-prompt side effects and 5 to 6.5 times less compute; (4) a control showing that rank 1 is a weak test (an additive edit forces random words to rank 1 on 20 of 20 prompts), which led to withdrawing an early ``knowledge absent'' classification; and (5) an applied side study with fixed pass bars: a linear translator between GPT-2 Small and Medium, Export and Import tests with controls, a documented search that locates the limit, and a first positive Import result.

\section{Organisation of the Report}

Chapter~\ref{ch:literature-review} reviews related work, including where gradient-based optimisation of edits comes from. Chapter~\ref{ch:method} gives the mathematics of the final system, Chapter~\ref{ch:evolution} the story of how it was reached, and Chapter~\ref{ch:transfer} the cross-model side study. Chapter~\ref{ch:implementation} describes the application and engineering, and Chapter~\ref{ch:conclusion} the findings, limits and future work. The appendices map journal entries to sections, list the side study's pass bars and give reproduction commands.
